\documentclass[11pt]{article}
\usepackage[T1]{fontenc}
\usepackage[utf8]{inputenc}
\usepackage{lmodern}
\usepackage[margin=1in]{geometry}
\usepackage{amsmath,amssymb,amsthm,mathtools}
\usepackage[hidelinks]{hyperref}
\usepackage{xcolor}

\setlength{\parindent}{0pt}
\setlength{\parskip}{0.62em}
\definecolor{statusblue}{RGB}{24,73,120}
\newcommand{\C}{\mathbb C}
\newcommand{\F}{\mathcal F}
\newcommand{\Spec}{\operatorname{Spec}}
\newcommand{\diag}{\operatorname{diag}}

\title{Exact norms for the general quadratic Fock evolution operator\\
\large A literature-implied answer to Remark 4.6 of arXiv:1409.1262}
\author{Run \texttt{fa\_banach\_001} --- lane 10}
\date{12 August 2026}

\begin{document}
\maketitle

\textbf{Status.} \textcolor{statusblue}{\texttt{literature\_implied\_answer
(full exact characterization).}}
The open exact-norm question in Aleman--Viola is answered by conjugating the
operator to the standard Fock space and applying Neretin's 2011 norm theorem
for Gaussian integral operators.  The finite-dimensional identification
below does not appear in either source; the norm theorem itself predates the
question, so this is not claimed as a new theorem.

\section*{The source question}

Let
\[
 H_\Phi=\operatorname{Hol}(\C^n)\cap
 L^2(\C^n,e^{-2\Phi}\,dL),\qquad
 P=(Mz)\mathbin\cdot\partial_z,
\]
where \(\Phi\) is a strictly convex real quadratic form.  The evolution is
\[
  e^{-tP}u(z)=u(e^{-tM}z).
\]
Remark 4.6 of \cite{AV} says that the exact norm is known when the
pluriharmonic part of \(\Phi\) vanishes, while ``the question remains open in
the general case.''

Use the source paper's decomposition
\[
 \Phi(z)=\frac12|Gz|^2-\Re h(z),\qquad
 h(z)=\frac12z^T H z,
\]
where \(G\) is positive definite and \(H=H^T\).  Put
\[
 K=G^{-T}HG^{-1}.
\]
Strict convexity is equivalent to \(\|K\|<1\) (for example by Takagi
diagonalization).

\section*{Exact finite-dimensional answer}

Fix a time \(t\) for which \(e^{-tP}\) is bounded, and define
\[
 A_t=e^{-tM},\qquad B_t=GA_tG^{-1},\qquad
 C_t=K-B_t^T K B_t,
\]
\[
 S_t=\begin{pmatrix}C_t&B_t^T\\ B_t&0\end{pmatrix},
 \qquad J=\begin{pmatrix}-I_n&0\\0&I_n\end{pmatrix}.
\]
The matrix \(S_t\) is complex symmetric.  Let
\(s_1(t),\ldots,s_n(t)\in[0,1]\) be its Gaussian canonical invariants in
the sense of Neretin \cite[Chapter 6, Theorem 1.3]{Neretin}.  They are
finite-dimensional data: Neretin's Theorem 1.4 says that the roots of
\[
 \det\!\left((I_{2n}-S_t^*S_t)
 -\lambda(J-S_t^*JS_t)\right)=0                         \tag{1}
\]
come in pairs
\[
 \lambda=\pm\frac{1-s_j(t)^2}{1+s_j(t)^2},
 \qquad j=1,\ldots,n.                                  \tag{2}
\]
Thus the \(s_j(t)\) can be recovered from a generalized eigenvalue problem
of size \(2n\).

Neretin's exact norm formula gives
\[
 \boxed{
 \|e^{-tP}\|
 =\det(I_{2n}-S_t^*J S_tJ)^{-1/4}
   \prod_{j=1}^n(1+s_j(t)^2)^{1/2}.}                   \tag{3}
\]
The positive value of the fourth root is understood.  Formula (3) applies to
\emph{every bounded time}, including boundary/noncompact cases
\cite[Chapter 6, Theorem 2.2(b)]{Neretin}.  Whenever \(\|S_t\|<1\), the
equivalent formula
\[
 \boxed{
 \|e^{-tP}\|
 =\det(I_{2n}-S_t^*S_t)^{-1/4}
   \prod_{j=1}^n(1-s_j(t)^2)^{1/2}}                    \tag{4}
\]
is also available.  Equations (1)--(4) are an exact, computable
characterization involving only the matrices \(G,H,M,t\).

\section*{Why the Gaussian theorem applies}

The unitary map used in \cite[Lemma 2.3]{AV} is
\[
 Uf(z)=|\det G|f(Gz)e^{-h(z)}:H_{|z|^2/2}\longrightarrow H_\Phi.
\]
A direct substitution, with \(w=Gz\), gives
\begin{align*}
 U^*e^{-tP}Uf(w)
 &=e^{h(G^{-1}w)-h(A_tG^{-1}w)}f(GA_tG^{-1}w)\\
 &=\exp\!\left(\frac12w^TC_tw\right)f(B_tw).           \tag{5}
\end{align*}
After the harmless scalar normalization of the standard Fock measure, its
reproducing formula is
\[
 f(\zeta)=\int_{\C^n}e^{\zeta^T\bar u}f(u)e^{-|u|^2}\,d\mu(u).
\]
Consequently the operator in (5) has kernel
\[
 \exp\!\left\{\frac12w^TC_tw+w^TB_t^T\bar u\right\}
 =\exp\!\left\{\frac12
 \begin{pmatrix}w^T&\bar u^T\end{pmatrix}
 S_t\begin{pmatrix}w\\\bar u\end{pmatrix}\right\}. \tag{6}
\]
This is exactly Neretin's Gaussian operator \(\mathcal B[S_t]\), as defined
in \cite[Chapter 5, equation (1.4)]{Neretin}.  Because \(U\) is unitary,
\(\|e^{-tP}\|=\|\mathcal B[S_t]\|\), and (3)--(4) follow from
\cite[Chapter 6, Theorems 1.4 and 2.2]{Neretin}.

The two extremal checks are immediate.  If \(h=0\), then \(C_t=0\) and the
formula reduces to norm one whenever the pure composition operator is
bounded, agreeing with the source's orthogonal case.  If \(B_t=0\), (5) is
rank one and (4) reduces to
\(\det(I-C_t^*C_t)^{-1/4}\), the direct Gaussian-vector norm.

\section*{Stable large-time sharpening}

Assume now that \(\Spec M\subset\{\Re\lambda>0\}\).  Then \(B_t\to0\),
\(C_t\to K\), and (5) converges to the rank-one operator
\[
 T_\infty f(w)=f(0)\exp\!\left(\frac12w^TKw\right).
\]
The vacuum evaluation \(f\mapsto f(0)\) has norm one in the normalized Fock
space, while Neretin's Gaussian-vector identity gives
\[
 \left\|\exp\!\left(\frac12w^TKw\right)\right\|
 =\det(I-K^*K)^{-1/4}.
\]
The operator-norm return estimate in \cite[Proposition 4.4 and Remark 4.6]{AV}
therefore yields the sharper asymptotic
\[
 \boxed{
 \|e^{-tP}\|=\det(I-K^*K)^{-1/4}
 +O\!\left(t^{R(M)-1}e^{-t\rho(M)}\right).}             \tag{7}
\]
In particular,
\[
 \lim_{t\to\infty}\|e^{-tP}\|=\det(I-K^*K)^{-1/4}.
\]
Since \(\|K\|<1\), this limit equals one exactly when \(K=0\), equivalently
when \(h=0\).  Thus (7) supplies the missing value behind the source's
``if and only if'' statement.

\section*{Scope and search record}

This answers the exact operator-norm question for every bounded time and for
every strictly convex quadratic Fock weight in the source setup.  It does not
claim a simpler radical expression for the roots of the \(2n\)-dimensional
pencil (1), nor a new closed formula for deciding the source paper's already
characterized set of bounded times.

All four cheap run indexes were searched; no prior packet was found.  Queries
included 1409.1262, the paper title, exact norms, and Gaussian operators.  A
current web search on 12 August 2026
found later work on boundedness, compactness, essential norms, and special
weighted composition operators, but no paper explicitly connecting this
remark to Neretin's theorem.  The source bibliography contains no Neretin,
Olshanski, Howe, or Gaussian-operator reference.

\begin{thebibliography}{9}
\bibitem{AV}
A. Aleman and J. Viola,
\emph{On weak and strong solution operators for evolution equations coming
from quadratic operators},
J. Spectr. Theory 8 (2018), 33--121;
\href{https://arxiv.org/abs/1409.1262}{arXiv:1409.1262}.

\bibitem{Neretin}
Y. A. Neretin,
\emph{Lectures on Gaussian Integral Operators and Classical Groups},
EMS Series of Lectures in Mathematics, European Mathematical Society, 2011;
especially Chapters 5--6.
\end{thebibliography}

\end{document}
